% =====================================================================
%  #betterposter starter — betterportraitposter (GPL-3.0)
%
%  Every workaround below is here because the class fails SILENTLY without
%  it. Numbered references point at references/class-traps.md.
%
%  Build:  make pdf        Verify:  make check
% =====================================================================

% [trap 13] The class does a bare \RequirePackage{xcolor}, so passing options
% with \usepackage[dvipsnames]{xcolor} later is an option clash. Options must
% arrive BEFORE \documentclass.
\PassOptionsToPackage{dvipsnames}{xcolor}
\documentclass[a0paper,fleqn]{betterportraitposter}

% --- Fonts -----------------------------------------------------------
% [trap 1] The class loads cmbright, then lato[default] (which pulls in
% fontspec => TU encoding), then [T1]{fontenc}. Under XeTeX that leaves
% TU/cmbr undefined and everything falls back to Latin Modern Sans, with
% only a warning. The compile still exits 0. Re-assert the font explicitly.
%
% A CLEAN COMPILE DOES NOT PROVE THE FONTS RESOLVED — `make check-fonts` does.
\usepackage{fontspec}
\setsansfont{Lato}[
  Extension      = .ttf,
  UprightFont    = *-Regular,
  ItalicFont     = *-Italic,
  BoldFont       = *-Bold,
  BoldItalicFont = *-BoldItalic,
]
\renewcommand{\familydefault}{\sfdefault}

% Weight contrast is cheaper than size contrast. Lato ships nine weights, so
% the main finding can set Light against Black — a three-step jump that makes
% the connective text recede and the emphasised terms carry the weight.
%
% [trap 2] Ligatures=TeX is NOT inherited by \newfontfamily. Without it,
% apostrophes print as straight typewriter quotes and -- stays two hyphens.
\newfontfamily\displayfont{Lato}[
  Extension   = .ttf,
  UprightFont = *-Light,
  BoldFont    = *-Black,
  Ligatures   = TeX,
]

\usepackage[none]{hyphenat}   % no hyphenation in display type
\usepackage{tikz}
\usetikzlibrary{arrows.meta,positioning}

% [trap 12] The class defines \qrcode with a different signature (it places
% two ready-made images). Park it, then hand the name to the package, which
% generates a code from a URL at build time.
% \let ...\relax suffices: LaTeX's \@ifdefinable permits redefining \relax.
\let\qrcodebox\qrcode
\let\qrcode\relax
\usepackage{qrcode}

% --- Palette ---------------------------------------------------------
% MEASURE the contrast before committing; a colour's name guarantees nothing.
% PineGreen is a BRIGHT green: 2.47:1 against white, below even the 3:1
% large-text floor. Darkening keeps the hue identity.
%     python3 scripts/contrast.py --mix 'PineGreen!65!black'
% and then, after the first build, check what actually rendered:
%     python3 scripts/contrast.py --sample main.pdf
\colorlet{postercolor}{PineGreen!65!black}

% [class] all four already exist, so \renewcommand — \newcommand errors.
\renewcommand{\maincolumnbackgroundcolor}{postercolor}
\renewcommand{\maincolumnfontcolor}{white}
\renewcommand{\columnbackgroundcolor}{white}
\renewcommand{\columnfontcolor}{black}

% --- Metrics ---------------------------------------------------------
% [trap 10] \bottomboxheight below the class's 0.1 default. With margin=0in
% the page finishes ~1 mm short of full and the residual glue spills onto a
% second, INKLESS page. Shrink the box; never add negative space at the end,
% because \bottombox starts with \vfill and will simply stretch further,
% pushing the footer off the page edge.
\setlength{\mainfindingheight}{0.24\paperheight}
\setlength{\bottomboxheight}{0.085\paperheight}

% [trap 3] \fontsizemain ships as 116/220 — a 190% leading that puts ~70 mm
% between two display lines and pushes the first one off the top of the page.
\renewcommand{\fontsizemain}{\fontsize{116}{132}\selectfont}
% Keep the byline below the title hierarchy: 48 pt is \fontsizesection.
\renewcommand{\fontsizeauthor}{\fontsize{42}{52}\selectfont}

% [trap 11] \section prepends \vspace{2em} = 34 mm at 48 pt. Four headings
% would spend 136 mm a figure-heavy poster does not have.
% [trap 8] \leavevmode so the heading is a box, not a whatsit.
\renewcommand{\section}[1]{%
  \vspace{0.7em}{\fontsizesection\selectfont\textbf{\leavevmode #1}}\\[0.25em]}

% --- Helpers ---------------------------------------------------------
% Uniform figure spacing. \begin{center} is a trivlist and adds \topsep both
% above AND below; a bare tikzpicture adds nothing. Normalise every wrapper to
% {\centering ...\par} and control the gap from here alone.
\newcommand{\figgap}{\par\vspace{\baselineskip}}
\newcommand{\figcap}[2][0.3em]{\par\vspace{#1}%
  {\fontsize{26}{31}\selectfont\color{black!62}\raggedright #2\par}\vspace{0.5em}}

% [trap 12] \color{black} and the white background are the difference between
% a working code and an empty rectangle: \bottombox sets \color{white}, so the
% package would draw white modules on white while the log still reports success.
\newcommand{\posterqr}[2]{%
  \begin{minipage}[c]{0.13\textwidth}%
    {\setlength{\fboxsep}{5mm}\colorbox{white}{%
       \color{black}\qrcode[height=0.6\bottomboxheight]{#1}}}%
  \end{minipage}%
  \begin{minipage}[c]{0.16\textwidth}%
    \fontsize{34}{40}\selectfont\raggedright #2%
  \end{minipage}%
}

% [trap 8] minipage[t] is a \vtop and takes its height from the FIRST BOX; a
% leading \color is a whatsit, so the column silently drops one line.
\newcommand{\authorcol}[4]{%
  \begin{minipage}[t]{#1}\centering
    \leavevmode\color{#4}#2\par\vspace{0.28em}
    {\fontsize{34}{42}\selectfont#3\par}%
  \end{minipage}}

\begin{document}

% =====================================================================
%  MAIN FINDING — one sentence, jargon-free, readable at 3 m.
%  Someone reading ONLY the bold words must get the finding.
% =====================================================================
% The class wraps \mainfinding's argument in tabular{p{0.9\textwidth}}, so \\
% there is a ROW break, not a line break. An inner minipage restores normal
% paragraph behaviour and is easier to control.
\mainfinding{%
  \begin{minipage}{\linewidth}
    \centering\displayfont
    % [trap 4] the \par lives INSIDE the size group: \baselineskip is read at
    % \par time, so closing the group first would set these lines on the
    % OUTER leading.
    {\fontsize{116}{128}\selectfont
     \textbf{SIX WORDS THAT}\\
     % Hanging punctuation keeps two centred lines optically aligned.
     \textbf{STOP A PASSER-BY}\rlap{\textbf{.}}\par}
    \vspace{0.7em}
    {\fontsize{64}{92}\selectfont
     One sentence saying what you found. Bold the \textbf{two or three terms}
     that carry it, and bold \textbf{both sides} of any contrast --- one-sided
     contrast does not read at a glance.\par}
  \end{minipage}%
}

% =====================================================================
%  TITLE
% =====================================================================
\titlebox{%
  {\centering
   {\fontsizetitle\selectfont\textbf{Full Paper Title, Exactly As Published}\par}
   \vspace{0.6em}
   {\fontsizeauthor
    % Underlining the presenting author is the conference convention, so it
    % carries meaning. An icon must be interpreted and drags an icon font in.
    \authorcol{0.30\linewidth}{\underline{First Author}}{Institution}{\institutefontcolor}%
    \hspace{0.02\linewidth}%
    \authorcol{0.30\linewidth}{Second Author}{Institution}{\institutefontcolor}%
    \hspace{0.02\linewidth}%
    \authorcol{0.30\linewidth}{Third Author}{Institution}{\institutefontcolor}%
    \par}
   \par}%
}

% =====================================================================
%  BODY — 3 short blocks, then a full-width results band.
%  Total visible prose across the whole poster: 150-250 words.
% =====================================================================
\centerbox{%
  \begin{multicols}{3}

    \section{Where}
    Roughly 35 words. Say what the data is and where it came from. Cutting a
    block entirely beats compressing all three.
    \figgap
    % Every figure wrapper opens the same way — no \begin{center}.
    {\centering\begin{tikzpicture}[x=1mm,y=1mm]
       \fill[postercolor!20] (0,0) rectangle (60,40);
       \node at (30,20) {\fontsize{26}{26}\selectfont figure};
     \end{tikzpicture}\par}
    \figcap{A caption gives what neither the figure nor the text gives alone.}

    \columnbreak

    \section{When}
    Roughly 35 words. The method, in plain language. No equations --- they do
    not survive being read at two metres.
    \figgap
    {\centering\begin{tikzpicture}[x=1mm,y=1mm]
       \fill[postercolor!40] (0,0) rectangle (60,40);
       \node at (30,20) {\fontsize{26}{26}\selectfont figure};
     \end{tikzpicture}\par}
    \figcap{For a method figure, the best caption is an analogy.}

    \columnbreak

    \section{How}
    Roughly 30 words. Inputs and output, stated explicitly. Name the baselines
    you compare against.
    \figgap
    {\centering\begin{tikzpicture}[x=1mm,y=1mm]
       \fill[postercolor!60] (0,0) rectangle (60,40);
       \node[white] at (30,20) {\fontsize{26}{26}\selectfont figure};
     \end{tikzpicture}\par}
    % No math mode and no \texttt here on purpose: either one pulls an
    % undeclared family into the PDF (cmbright math, Latin Modern Mono) and
    % `make check-fonts` will correctly flag it as a fallback.
    \figcap{Rasters need 120 ppi or better at the size they are placed.}

  \end{multicols}

  % [trap 6] The full-width band lives INSIDE this same \centerbox, after
  % \end{multicols}. A second \centerbox would draw a second horizontal rule.
  \section{So what}
  % [trap 7] [b] aligns LAST baselines. [t] aligns FIRST baselines, which for
  % two different content types leaves one sitting tens of mm below the other
  % and inflates the whole band.
  \begin{minipage}[b]{0.63\textwidth}
    {\centering\begin{tikzpicture}[x=1mm,y=1mm]
       \fill[postercolor!15] (0,0) rectangle (140,70);
       \node at (70,35) {\fontsize{30}{30}\selectfont results figure};
     \end{tikzpicture}\par}
  \end{minipage}\hfill
  \begin{minipage}[b]{0.30\textwidth}
    {\centering\begin{tikzpicture}[x=1mm,y=1mm]
       \fill[postercolor!15] (0,0) rectangle (65,65);
       \node at (32,32) {\fontsize{30}{30}\selectfont second view};
     \end{tikzpicture}\par}
  \end{minipage}
}

% =====================================================================
%  FOOTER
% =====================================================================
% [trap 5] \bottombox injects space on BOTH sides of its argument — 8.4 mm at
% \fontsizesection's 48 pt. With \hfill already spending the slack, the content
% overflows the right margin. The 0.97 minipage absorbs it.
\bottombox{%
  \begin{minipage}{0.97\linewidth}
    \posterqr{https://example.org/paper}{Paper}%
    \hfill
    \posterqr{https://example.org/code}{Code}%
    \hfill
    % A dark logo vanishes on the coloured band — put logos on a white panel,
    % which also matches the light ground the QR codes already need.
    \begin{minipage}[c]{0.22\textwidth}\raggedleft
      {\setlength{\fboxsep}{4mm}\colorbox{white}{%
         % \includegraphics[height=0.5\bottomboxheight]{assets/logo.pdf}
         \rule{0pt}{0.5\bottomboxheight}\textcolor{black}{\ \ logo\ \ }}}%
    \end{minipage}
  \end{minipage}%
}

\end{document}
